\chapter{1995年上海卷（文）}
\section{单选题}

\begin{problem}
\(y=\sin^2 x\) 是（\quad）

\choices
    {最小正周期为 \(2\pi \) 的偶函数}
    {最小正周期为 \(2 \pi\) 的奇函数}
    {最小正周期为 \(\pi \) 的偶函数}
    {最小正周期为 \(\pi\) 的奇函数}
\end{problem}

\begin{answer}
C
\end{answer}

\begin{problem}
如果，\(P = \{x \mid (x - 1)(2x - 5) < 0\}\)，\(Q = \{x \mid 0 < x < 10\}\)，那么（\quad）

\choices
    {\(P\cap Q=\emptyset\)}
    {\(P \subset Q\)}
    {\(P\supset Q\)}
    {\(P\cup Q=\R\)}
\end{problem}

\begin{answer}
B
\end{answer}

\begin{problem}
方程 \(\tan\left(2x+\frac{\pi}{3}\right)=\frac{\sqrt{3}}{3}\) 在区间 \([0,2\pi)\) 上解的个数是（\quad）

\choices
    {\(5\)}
    {\(4\)}
    {\(3\)}
    {\(2\)}
\end{problem}

\begin{answer}
B
\end{answer}

\begin{problem}
设棱锥的底面面积是 \(8\mathrm{~cm}^2\)，那么这个棱锥的中截面（过棱锥高的中点且平行于底面的截面）的面积是（\quad）

\choices
    {\(4\mathrm{~cm}^2\)}
    {\(2\sqrt{2}\mathrm{~cm}^2\)}
    {\(2\mathrm{~cm}^2\)}
    {\(\sqrt{2}\mathrm{~cm}^2\)}
\end{problem}

\begin{answer}
C
\end{answer}

\begin{problem}
“\(ab<0\)”是“方程 \(ax^2+by^2=c\) 表示双曲线”的（\quad）

\choices
    {必要条件但不是充分条件}
    {充分条件但不是必要条件}
    {充分必要条件}
    {既不是充分条件又不是必要条件}
\end{problem}

\begin{answer}
A
\end{answer}

\begin{problem}
当 \(a\neq 0\) 时，函数 \(y=ax+b\) 和 \(y=b^{ax}\) 的图象只可能是（\quad）

\choices
    {\choicebitmap[width=0.15\paperwidth]{img/1995/shanghai_liberal/q6_optA_nolabel.png}}
    {\choicebitmap[width=0.15\paperwidth]{img/1995/shanghai_liberal/q6_optB_nolabel.png}}
    {\choicebitmap[width=0.15\paperwidth]{img/1995/shanghai_liberal/q6_optC_nolabel.png}}
    {\choicebitmap[width=0.15\paperwidth]{img/1995/shanghai_liberal/q6_optD_nolabel.png}}
\end{problem}

\begin{answer}
A
\end{answer}

\begin{problem}
当 \(0<a<b<1\) 时，下列不等式中正确的是（\quad）

\choices
    {\((1 - a)^{\frac{1}{b}} > (1 - a)^{b}\)}
    {\((1 + a)^{a} > (1 + b)^{b}\)}
    {\((1 - a)^{b} > (1 - a)^{\frac{b}{2}}\)}
    {\((1 - a)^{a} > (1 - b)^{b}\)}
\end{problem}

\begin{answer}
D
\end{answer}

\begin{problem}
下列四个命题中的真命题是（\quad）

\choices
    {经过定点 \(P_0(x_0,y_0)\) 的直线都可以用方程 \(y-y_0=k(x-x_0)\) 表示}
    {经过任意两个不同的点 \(P_1(x_1,y_1)\)，\(P_2(x_2,y_2)\) 的直线都可以用方程 \((y-y_1)(x_2-x_1)=(x-x_1)(y_2-y_1)\) 表示}
    {不经过原点的直线都可以用方程 \(\frac{x}{a}+\frac{y}{b}=1\) 表示}
    {经过定点 \(A(0,b)\) 的直线都可以用方程 \(y=kx+b\) 表示}
\end{problem}

\begin{answer}
B
\end{answer}

\section{填空题}

\begin{problem}
不等式 \(\frac{2x-1}{x+3}>1 \) 的解是\fillinblank{}.
\end{problem}

\begin{answer}
\(x<-3\) 或 \(x>4\)
\end{answer}

\begin{problem}
双曲线 \(\frac{x^{2}}{2}-\frac{8y^{2}}{9}=8 \) 的渐近线方程是\fillinblank{}.
\end{problem}

\begin{answer}
\(3x\pm4y=0\)
\end{answer}

\begin{problem}
1992年底世界人口达到 \(54.8\text{ 亿}\)，若人口的年平均增长率为 \(x\%\)，2000 年底世界人口数为 \(y\)（\(\text{亿}\)），那么 \(y\) 与 \(x\) 的函数关系式是\fillinblank{}.
\end{problem}

\begin{answer}
\(y=54.8\left(1+x\%\right)^8\)
\end{answer}

\begin{problem}
已知 \(\log_3 x=-\frac{1}{\log_2 3}\)，那么 \(x+x^2+x^3+\cdots+x^n+\cdots=\)\fillinblank{}.
\end{problem}

\begin{answer}
1
\end{answer}

\begin{problem}
若 \((x+1)^n=x^n+\cdots+ax^3+bx^2+\cdots+1\)（\(n\in\N\)），且 \(a:b=3:1\)，那么 \(n=\)\fillinblank{}.
\end{problem}

\begin{answer}
11
\end{answer}

\begin{problem}
到点 \(A(-1,0)\) 和直线 \(x = 3\) 距离相等的点的轨迹方程是\fillinblank{}.
\end{problem}

\begin{answer}
\(y^2=-8x+8\)
\end{answer}

\begin{problem}
函数 \(y=3x^2+1\)（\(x\leqslant 0\)）的反函数是 \(y=\)\fillinblank{}.
\end{problem}

\begin{answer}
\(-\sqrt{\frac{x-1}{3}}\)（\(x\geqslant1\)）
\end{answer}

\begin{problem}
函数 \(y=\lg\sqrt{10^{x}-2} \) 的定义域是\fillinblank{}.
\end{problem}

\begin{answer}
\((\lg2,+\infty)\)
\end{answer}

\begin{problem}
函数 \(y=\sin\frac{x}{2}+\cos\frac{x}{2} \) 在 \((-2\pi,2\pi) \) 内的递增区间是\fillinblank{}.
\end{problem}

\begin{answer}
\(\left[-\frac{3\pi}{2},\frac{\pi}{2}\right]\)
\end{answer}

\section{选做题：填空题}

\begin{problem}[A]
若 \(\lim\limits_{n\to\infty}\left[1+(r+1)^n\right]=1\)，则 \(r\) 的取值范围是\fillinblank{}.
\end{problem}

\begin{answer}
\(-2<r<0\)
\end{answer}

\begin{problem}[A]
从\(6\text{ 台}\)原装计算机和\(5\text{ 台}\)组装计算机中任意选取\(5\text{ 台}\)，其中至少有原装与组装计算机各\(2\text{ 台}\)，则不同的选取法有\fillinblank{} 种（结果用数值表示）.
\end{problem}

\begin{answer}
350
\end{answer}

\begin{problem}[A]
把圆心角为 \(216^{\circ} \)，半径为 \(5\text{ 分米}\) 的扇形铁皮焊成一个圆锥形容器（不计焊缝），那么容器的容积是\fillinblank{} \(\text{立方分米}\)（结果保留两位小数）.
\end{problem}

\begin{answer}
\(37.70\)（或 \(37.68\) 或 \(37.69\)）
\end{answer}

\setcounter{probnum}{17}
\begin{problem}[B]
\(\lim\limits_{n\to\infty}\left(1+\frac{1}{n}\right)^{n-2}=\)\fillinblank{}.
\end{problem}

\begin{answer}
\(e\)
\end{answer}

\begin{problem}[B]
从\(6\text{ 台}\)原装计算机和\(5\text{ 台}\)组装计算机中任意选取\(5\text{ 台}\)，其中至少有原装与组装计算机各\(2\text{ 台}\)的概率是\fillinblank{}.
\end{problem}

\begin{answer}
\(\frac{25}{33}\)（或 \(0.76\)）
\end{answer}

\begin{problem}[B]
某工程由下列工序组成，那么工程总时数是\fillinblank{} \(\text{天}\).

\begin{center}
\begin{tabular}{|c|c|c|c|c|c|}
\hline
工序 & \(a\) & \(b\) & \(c\) & \(d\) & \(e\) \\
\hline
紧前工序 & — & \(a\) & \(a\) & \(b\) & \(c\)，\(d\) \\
\hline
工时数（天） & \(2\) & \(3\) & \(5\) & \(6\) & \(3\) \\
\hline
\end{tabular}
\end{center}
\end{problem}

\begin{answer}
14
\end{answer}

\section{选做题：解答题}

\begin{problem}[A]
已知 \(\tan\left(\frac{\pi}{4}+\theta\right)=3\)，求 \(\sin2\theta-2\cos^{2}\theta\) 的值.
\end{problem}

\begin{solution}
因为
\[
\frac{1+\tan\theta}{1-\tan\theta}=3,
\]
解得 \(\tan\theta=\frac{1}{2}\). 于是
\[
\sin2\theta-2\cos^2\theta
=\sin2\theta-\cos2\theta-1
=\frac{2\tan\theta}{1+\tan^2\theta}-\frac{1-\tan^2\theta}{1+\tan^2\theta}-1
=\frac{4}{5}-\frac{3}{5}-1=-\frac{4}{5}.
\]
\end{solution}

\begin{problem}[A]
已知复数 \(z_1\)，\(z_2\) 满足 \(|z_1|=|z_2|=1\)，且 \(z_1+z_2=\i\). 求 \(z_1\)，\(z_2\) 的值.
\end{problem}

\begin{solution}
因为 \(z_1+z_2=\i\) 是纯虚数，并且 \(|z_1|=|z_2|=1\)，可设
\[
z_1=a+b\i,\qquad z_2=-a+b\i\quad(a,b\in\R).
\]
且
\[
a^2+b^2=1.
\]
于是由
\[
(a+b\i)+(-a+b\i)=\i
\]
可得
\[
b=\frac{1}{2},\qquad a=\pm\frac{\sqrt{3}}{2}.
\]
所以
\[
z_1=\frac{\sqrt{3}}{2}+\frac{1}{2}\i,\qquad
z_2=-\frac{\sqrt{3}}{2}+\frac{1}{2}\i,
\]
或
\[
z_1=-\frac{\sqrt{3}}{2}+\frac{1}{2}\i,\qquad
z_2=\frac{\sqrt{3}}{2}+\frac{1}{2}\i.
\]
\end{solution}

\setcounter{probnum}{20}
\begin{problem}[B]
如图，在空间直角坐标系中，\(BC=2\)，原点 \(O\) 是 \(BC\) 的中点，点 \(A\) 的坐标是 \(\left(\frac{\sqrt{3}}{2},\frac{1}{2},0\right)\)，点 \(D\) 在平面 \(yoz\) 上，且 \(\angle BDC=90^\circ\)，\(\angle DCB=30^\circ\).

\begin{enumerate}
\item 求向量 \(\overrightarrow{OD}\) 的坐标；
\item 设向量 \(\overrightarrow{AD}\) 和 \(\overrightarrow{BC}\) 的夹角为 \(\theta\)，求 \(\cos\theta\) 的值.
\end{enumerate}

\FigureLayoutDeclare{img/1995/shanghai_liberal/q21_fig1.png}{6.6cm}{0bp 0bp 0bp 0bp}
\bitmapfigure[max width=0.42\linewidth]{img/1995/shanghai_liberal/q21_fig1.png}
\end{problem}

\begin{solution}
\begin{enumerate}
\item 因为 \(z_D=|BC|\cos30^\circ\sin30^\circ=\frac{\sqrt{3}}{2}\)，\(y_D=1-|BC|\cos30^\circ\cos30^\circ=-\frac{1}{2}\)，所以
\[
\overrightarrow{OD}=\left(0,-\frac{1}{2},\frac{\sqrt{3}}{2}\right).
\]

\item 因为 \(\overrightarrow{OA}=\left(\frac{\sqrt{3}}{2},\frac{1}{2},0\right)\)，\(\overrightarrow{OB}=(0,-1,0)\)，\(\overrightarrow{OC}=(0,1,0)\)，所以
\[
\overrightarrow{AD}=\left(-\frac{\sqrt{3}}{2},-1,\frac{\sqrt{3}}{2}\right),\qquad
\overrightarrow{BC}=(0,2,0).
\]
于是
\[
\cos\theta=\frac{\overrightarrow{AD}\cdot\overrightarrow{BC}}{|\overrightarrow{AD}||\overrightarrow{BC}|}
=\frac{-2}{\sqrt{\frac{5}{2}}\cdot2}=-\frac{\sqrt{10}}{5}.
\]
\end{enumerate}
\end{solution}

\begin{problem}[B]
设 \(y=f(x)\) 是二次函数，方程 \(f(x)=0\) 有两个相等的实根，且 \(f'(x)=2x+2\).

\begin{enumerate}
\item 求 \(y=f(x)\) 的表达式；
\item 求 \(y=f(x)\) 的图象与两坐标轴所围成图形的面积.
\end{enumerate}
\end{problem}

\begin{solution}
\begin{enumerate}
\item 设 \(f(x)=ax^2+bx+c\)，则 \(f'(x)=2ax+b\). 又因为 \(f'(x)=2x+2\)，所以 \(a=1\)，\(b=2\)，从而
\[
f(x)=x^2+2x+c.
\]
方程 \(f(x)=0\) 有两个相等的实根，所以判别式 \(\Delta=4-4c=0\)，即 \(c=1\). 因此
\[
f(x)=x^2+2x+1.
\]

\item 所求面积
\[
=\int_{-1}^{0}(x^2+2x+1)\mathrm{d}x
=\left(\frac{1}{3}x^3+x^2+x\right)\Big|_{-1}^{0}
=\frac{1}{3}.
\]
\end{enumerate}
\end{solution}

\section{解答题}

\begin{problem}
如图，四棱锥 \(P-ABCD\) 中，底面是一个矩形，\(AB=3\)，\(AD=1\)，又 \(PA\perp AB\)，\(PA=4\)，\(\angle PAD=60^\circ\).

\begin{enumerate}
\item 求四棱锥 \(P-ABCD\) 的体积；
\item 求二面角 \(P-BC-D\) 的大小（用反三角函数表示）.
\end{enumerate}

\bitmapfigure[max width=0.34\linewidth]{img/1995/shanghai_liberal/q23_fig1.png}
\end{problem}

\begin{solution}
\begin{enumerate}
\item 因为 \(AB\perp AD\)，\(AB\perp AP\)，所以 \(AB\perp\) 平面 \(PAD\). 又 \(AB\subset\) 平面 \(ABCD\)，所以平面 \(ABCD\perp\) 平面 \(PAD\).

在平面 \(PAD\) 中，作 \(PE\perp AD\)（\(AD=\) 平面 \(PAD\cap\) 平面 \(ABCD\)），交 \(AD\) 的延长线于点 \(E\)（因为 \(AE=AP\cos60^\circ=2>AD\)）. 所以 \(PE\perp\) 平面 \(ABCD\).

在直角三角形 \(PAE\) 中，\(PE=AP\sin60^\circ=2\sqrt{3}\). 因此
\[
V_{P-ABCD}=\frac{1}{3}AB\cdot AD\cdot PE=2\sqrt{3}.
\]

\item 在平面 \(ABCD\) 中，作 \(EF\parallel DC\)，交 \(BC\) 的延长线于点 \(F\)，则 \(EF\perp BF\)，连接 \(PF\). 因为 \(PE\perp\) 平面 \(ABCD\)，且 \(EF\perp BF\)，所以 \(PF\perp BF\)（三垂线定理）.

于是 \(\angle PFE\) 是二面角 \(P-BC-D\) 的平面角. 在直角三角形 \(PEF\) 中，
\[
\tan\angle PFE=\frac{PE}{EF}=\frac{2\sqrt{3}}{3}.
\]
因此，二面角 \(P-BC-D\) 的大小为
\[
\arctan\frac{2\sqrt{3}}{3}.
\]
\end{enumerate}
\end{solution}

\begin{problem}
设椭圆的方程为 \(\frac{x^2}{m^2}+\frac{y^2}{n^2}=1\)（\(m,n>0\)），过原点且倾角为 \(\theta\) 和 \(\pi-\theta\)（\(0<\theta<\frac{\pi}{2}\)）的两条直线分别交椭圆于 \(A\)，\(C\) 和 \(B\)，\(D\) 四点.

\begin{enumerate}
\item 用 \(\theta\)，\(m\)，\(n\) 表示四边形 \(ABCD\) 的面积 \(S\)；
\item 若 \(m\)，\(n\) 为定值，当 \(\theta\) 在 \(\left(0,\frac{\pi}{4}\right]\) 上变化时，求 \(S\) 的最大值 \(u\).
\end{enumerate}
\end{problem}

\begin{solution}
\begin{enumerate}
\item 设经过原点且倾角为 \(\theta\) 的直线方程为 \(y=x\tan\theta\)，则方程组为
\[
\begin{cases}
y=x\tan\theta,\\
\dfrac{x^2}{m^2}+\dfrac{y^2}{n^2}=1.
\end{cases}
\]
解得
\[
\begin{cases}
x^2=\dfrac{m^2n^2}{n^2+m^2\tan^2\theta},\\
y^2=\dfrac{m^2n^2\tan^2\theta}{n^2+m^2\tan^2\theta}.
\end{cases}
\]
由对称性，得四边形 \(ABCD\) 为矩形，又因为 \(0<\theta<\frac{\pi}{2}\)，所以四边形 \(ABCD\) 的面积
\[
S=4|x||y|=\frac{4m^2n^2\tan\theta}{n^2+m^2\tan^2\theta}.
\]

\item
\[
S=\frac{4m^2n^2}{\dfrac{n^2}{\tan\theta}+m^2\tan\theta}.
\]

\begin{enumerate}
\item 当 \(m>n\)，即 \(\frac{n}{m}<1\) 时，
\[
\frac{n^2}{\tan\theta}+m^2\tan\theta\geqslant2mn,
\]
且仅当 \(\tan^2\theta=\frac{n^2}{m^2}\) 时等号成立. 因此
\[
S\leqslant\frac{4m^2n^2}{2mn}=2mn.
\]
由于 \(0<\theta\leqslant\frac{\pi}{4}\)，\(0<\tan\theta\leqslant1\)，故取 \(\tan\theta=\frac{n}{m}\)，得 \(u=2mn\).

\item 当 \(m<n\)，即 \(\frac{n}{m}>1\) 时，对于任意 \(0<\theta_1<\theta_2\leqslant\frac{\pi}{4}\)，由于
\[
\begin{aligned}
&\left(m^2\tan\theta_2+\frac{n^2}{\tan\theta_2}\right)-\left(m^2\tan\theta_1+\frac{n^2}{\tan\theta_1}\right)\\
&=(\tan\theta_2-\tan\theta_1)\frac{m^2\tan\theta_1\tan\theta_2-n^2}{\tan\theta_1\tan\theta_2}<0,
\end{aligned}
\]
所以 \(S\) 是 \(\theta\) 的递增函数，取 \(\theta=\frac{\pi}{4}\)，得
\[
u=\frac{4m^2n^2}{m^2+n^2}.
\]
\end{enumerate}
\end{enumerate}
\end{solution}

\begin{problem}
已知二次函数 \(y=f(x)\) 在 \(x=\frac{t+2}{2}\) 处取得最小值 \(-\frac{t^2}{4}\)（\(t\neq 0\)），\(f(1)=0\).

\begin{enumerate}
\item 求 \(y=f(x)\) 的表达式；
\item 若任意实数 \(x\) 都满足等式 \(f(x)g(x)+mx+n=x^3\)（\(g(x)\) 为多项式），试用 \(t\) 表示 \(m\) 和 \(n\)；
\item 设 \(t=-3\)，圆 \(C_k\) 的方程为 \((x-mq^k)^2+(y-nq^k)^2=r_k^2\)（\(0<q<1\)，\(r_k>0\)），圆 \(C_k\) 与 \(x\) 轴相切，圆 \(C_k\) 与 \(C_{k+1}\) 外切（\(k=0\)，\(1\)，\(2\)，\(\cdots\)），记 \(S\) 为所有这些圆的面积之和，求 \(q\)，\(S\) 的值.
\end{enumerate}
\end{problem}

\begin{solution}
\begin{enumerate}
\item 设
\[
f(x)=a\left(x-\frac{t+2}{2}\right)^2-\frac{t^2}{4}.
\]
因为 \(f(1)=0\)，所以
\[
a\left(1-\frac{t+2}{2}\right)^2-\frac{t^2}{4}=0,
\]
从而 \(a=1\). 于是
\[
f(x)=x^2-(t+2)x+t+1.
\]

\item 因为 \(f(x)=(x-1)[x-(t+1)]\)，代入已知等式，得
\[
(x-1)[x-(t+1)]g(x)+mx+n=x^3.
\]
将 \(x=1\)，\(x=t+1\) 分别代入上式，得
\[
\begin{cases}
m+n=1,\\
(t+1)m+n=(t+1)^3.
\end{cases}
\]
由 \(t\neq0\)，可解得
\[
m=t^2+3t+3,\qquad n=-t^2-3t-2.
\]

\item 当 \(t=-3\) 时，\(m=3\)，\(n=-2\). 圆 \(C_k\) 的方程为
\[
(x-3q^k)^2+(y+2q^k)^2=r_k^2.
\]
因为圆 \(C_k\) 与 \(x\) 轴相切，所以 \(r_k=2q^k\).

圆 \(C_k\) 的圆心 \(O_k=(3q^k,-2q^k)\)，所以
\[
|O_kO_{k+1}|^2=(3q^{k+1}-3q^k)^2+(2q^{k+1}-2q^k)^2
=13q^{2k}(q-1)^2.
\]
因为 \(0<q<1\)，所以
\[
|O_kO_{k+1}|=\sqrt{13}q^k(1-q).
\]
又圆 \(C_k\) 与 \(C_{k+1}\) 外切，所以
\[
r_k+r_{k+1}=|O_kO_{k+1}|,
\]
即
\[
2q^k+2q^{k+1}=\sqrt{13}q^k(1-q).
\]
解得
\[
q=\frac{\sqrt{13}-2}{\sqrt{13}+2}.
\]

圆 \(C_k\) 的面积
\[
S_k=\pi r_k^2=4\pi q^{2k},
\]
所以
\[
S=S_0+S_1+S_2+\cdots
=4\pi(1+q^2+q^4+\cdots)
=\frac{4\pi}{1-q^2}
=\pi\left(2+\frac{17}{26}\sqrt{13}\right).
\]
\end{enumerate}
\end{solution}
